\documentclass[10pt,a4paper]{article}
\usepackage[margin=0.6in]{geometry}
\usepackage{amsmath,amssymb,amsfonts,stmaryrd}
\usepackage{xcolor}
\usepackage{enumitem}
\usepackage{tcolorbox}
\usepackage{multicol}
\usepackage{listings}
\usepackage{hyperref}

\definecolor{isabelleblue}{RGB}{0, 70, 140}
\definecolor{isabellekw}{RGB}{120, 0, 120}
\definecolor{boxbg}{RGB}{248, 249, 250}

\tcbset{
  colback=boxbg,
  colframe=isabelleblue,
  fonttitle=\bfseries\large,
  boxrule=0.8pt,
  arc=3pt,
  left=8pt,
  right=8pt,
  top=6pt,
  bottom=6pt
}

\title{\textbf{\huge Forward Reasoning in Isabelle (FW)}\\[0.3em]
\large Lecture Notes}
\author{\textbf{Course:} Computer-Assisted Theorem Proving (7CCSACTP)}
\date{}

\begin{document}

\maketitle
\vspace{-1.5em}

\begin{tcolorbox}[title={1. Context \& Overview}]
\begin{multicols}{2}
\textbf{$\rightarrow$ \underline{Last Video Revision}:}
\begin{itemize}[leftmargin=1.2em, itemsep=2pt]
    \item \textbf{Basic Functional Programming:} Currying, Lambda expressions, Polymorphic types.
    \item \textbf{Isabelle Interface:} Theory files, terms (numbers, lists, sets, products/tuples, functions) and type annotations.
    \item \textbf{Logical Specifications:} Connectives and theorem structure.
\end{itemize}

\columnbreak

\textbf{$\rightarrow$ \underline{Performing Proofs in Isabelle}:}
\begin{itemize}[leftmargin=1.2em, itemsep=2pt]
    \item \textbf{Logical Reasoning:} Forward reasoning (FW) \& Backward reasoning (BW).
    \item \textbf{Proofs by Substitution.}
    \item \textbf{Mathematical Induction} on natural numbers (\texttt{nat}).
\end{itemize}
\end{multicols}
\end{tcolorbox}

\vspace{0.3em}

\begin{tcolorbox}[title={2. Forward Reasoning Mechanics: Example Proof}]
\textbf{Goal Statement:} $P \rightarrow Q, \quad Q \rightarrow R, \quad \neg R \;\vdash\; \neg P$

\begin{multicols}{2}
\textbf{Step-by-Step Derivation:}
\begin{enumerate}[leftmargin=1.5em, itemsep=3pt]
    \item $\neg R$ \quad \textit{(Assumption 1)}
    \item $Q \rightarrow R$ \quad \textit{(Assumption 2)}
    \item $Q \rightarrow \neg R$ \quad \textit{(from 1, using $\rightarrow I$ rule)}
    \item $\neg Q$ \quad \textit{(from 2, 3, using $\neg I$ rule)}
    \item $P \rightarrow \neg Q$ \quad \textit{(from 4, using $\rightarrow I$ rule)}
    \item $\neg P$ \quad \textit{(from assumption $P \rightarrow Q$ and 5, using $\neg I$)} $\checkmark$
\end{enumerate}

\columnbreak

\textbf{Inference Rules \& Isabelle Syntax:}
\begin{itemize}[leftmargin=1.2em, itemsep=3pt]
    \item Implication Introduction ($\rightarrow I$):
    \[ \frac{X}{Y \Rightarrow X} \]
    \item Negation Introduction ($\neg I$):
    \[ \frac{X \rightarrow Y \qquad X \rightarrow \neg Y}{\neg X} \]
    \item \textbf{Isabelle Forward Attribute (\texttt{OF}):}
    \begin{center}
    \texttt{$f_2$ [OF $f_1$]}
    \end{center}
    Instantiates premise of rule $f_2$ with known fact $f_1$ (e.g. \texttt{impI[OF fact1]}).
\end{itemize}
\end{multicols}
\end{tcolorbox}

\vspace{0.3em}

\begin{tcolorbox}[title={3. Isabelle Forward Reasoning Notes}]
\begin{itemize}[leftmargin=1.2em, itemsep=3pt]
    \item \textbf{Using Facts:} To obtain and utilize a fact $P$ in Isar:
    \begin{center}
    \texttt{have P \quad using P}
    \end{center}
    \item \textbf{Schematic Variables:} In Isabelle, schematic variables written with a question mark (e.g. \texttt{?x}) represent universally quantified variables ($\forall x$).
\end{itemize}
\end{tcolorbox}

\end{document}
